\chapter{Conclusions}\label{C:con}
This report investigated how machine learning can be used to make
RPL's parent selection load-aware. The objective functions proposed by
the ROLL working group rely on a single metric and react to
congestion only after a route has degraded. To address this, we
proposed MLOF, an objective function that uses a
lightweight regression model to predict the packet delivery ratio
of a candidate parent's path. The predicted PDR is combined with ETX
to compute the rank increase. To train the model, we collected data from
random Cooja topologies containing overloading clients and compared
several regression models. We then used Spearman correlation,
permutation importance and SHAP to reduce the input to six features:
\texttt{is\_new}, \texttt{cpu}, \texttt{path\_cpu},
\texttt{drop\_rate}, \texttt{path\_ppm} and \texttt{hop\_count}. The
feature analysis showed that the metrics describing the load on the
path towards the root, such as the parent's drop rate, CPU usage and
traffic rate, together with the hop count, are far more informative of
the delivery ratio than link-level metrics such as RSSI and ETX. Two
models, a decision tree (DTree) and a gradient-boosted ensemble
(LGBM), were embedded in Contiki-NG on the Z1 platform. We evaluated
them against OF0 and MRHOF in random topologies with 60, 80 and 100
nodes and traffic loads from 64 to 128~bps per node.

The evaluation shows that MLOF, and DTree in particular, extends the
range of load that the network can sustain. At 60 nodes and 128~bps,
DTree delivered 94.2\% of packets, compared with 83.1\% for OF0 and
76.8\% for MRHOF. At 100 nodes and 64~bps it delivered 97.6\%, against
81.5\% for OF0. DTree also formed the shortest paths over
the best links, which resulted in the lowest end-to-end latency and
CPU usage in every configuration. Furthermore, MLOF built a
considerably more stable DODAG. It performed a fraction of the parent
switches of OF0 and MRHOF and sent fewer DIO messages, so that its
total routing control traffic remained lower despite the 7 bytes added
to
each DIO. Inference on the Z1 took 13.0~$\mu$s for DTree and
118.8~$\mu$s for LGBM, which shows that the cost of embedding a
tree-based model in the objective function is small compared with the
savings in routing overhead.

However, the advantage of MLOF is limited to the range before the
network saturates. Once the load exceeds the capacity of the network,
at 80 nodes from 112~bps and at 100 nodes from 96~bps, all OFs suffer
a similar loss of delivery ratio and DTree performs no better than
OF0.
In this regime, LGBM also maintained a more stable topology than
DTree,
although the simpler decision tree performed best overall. These
results indicate that a load-aware objective function can delay
congestion but cannot compensate for insufficient network capacity.

\section{Limitations}

\begin{itemize}
    \item \textbf{No radio duty cycling.} All simulations used CSMA
        with the radio always on. Radio listening therefore dominates
        the energy consumption, so the reduction in CPU usage and
        control traffic could not be translated into a meaningful
        evaluation of energy consumption or network lifetime.
    \item \textbf{Simulation only.} The model was both trained and
        evaluated on data generated in Cooja. Although Cooja emulates
        the Z1 hardware at the instruction level, its radio model does
        not capture the interference, multipath fading and
        time-varying link quality of real deployments. The model may
        therefore need to be retrained or calibrated before it can be
        deployed on physical nodes.
    \item \textbf{Limited range of scenarios.} The training and
        evaluation topologies were drawn from the same random
        generator, with a single root, static nodes and a fixed share
        of 5\% overloading clients. It remains open whether the model
        generalises to other layouts, such as grid, linear or
        clustered topologies, and to mobile nodes, multiple roots or
        different traffic patterns.
\end{itemize}

\section{Future Work}

\begin{itemize}
    \item Evaluate MLOF with radio duty cycling, for example
        ContikiMAC or TSCH, to quantify its effect on energy
        consumption and network lifetime. In duty-cycled networks,
        transmissions and processing make up a larger share of the
        energy budget, so the gains in stability are expected to be
        more pronounced.
    \item Validate MLOF on a physical testbed, such as FIT IoT-LAB, to
        assess whether a model trained in simulation transfers to
        real radio conditions, and investigate retraining on
        testbed data.
    \item Extend the evaluation to a wider range of topologies and
        traffic patterns, including different shares and placements of
        overloading clients, and assess how well the model
        generalises to scenarios that were not seen during training.
    \item Improve the behaviour of MLOF under saturation, for example
        by tuning the weight $\alpha$ between the predicted PDR and
        ETX, by adding a stronger switching hysteresis, or by
        combining
        it with admission or rate control at the application layer.
    \item Encode the shared metrics according to the RFC~6551 metric
        container format to improve interoperability with standard
        RPL implementations.
\end{itemize}
